\documentclass[12pt]{letter}
\usepackage[margin=1in]{geometry}
\signature{Showerhead}
\address{}
\begin{document}
\begin{letter}{Editor-in-Chief \\ \textit{Cryptologia}}
\opening{Dear Editor,}

We submit for your consideration the manuscript ``The Voynich Manuscript as Lossy Generation,'' which presents quantitative evidence that the Voynich Manuscript's text was produced by a physical combinatorial generation device.

This work is particularly suited to \textit{Cryptologia} for three reasons. First, it directly engages and computationally falsifies the Cardan grille hypothesis of Rugg (2004), published in this journal. Second, it extends the generation model of Timm \& Schinner (2020), also published here. Third, it proposes a new analytical framework---lossy generation---that provides a middle ground between the cipher and hoax hypotheses that have dominated the Voynich literature.

The paper presents five convergent lines of evidence: positional slot dependencies ruling out independent-column mechanisms, a default-filler signature diagnostic of table lookup, independent recovery of manuscript section structure, a synthetic generator comparison falsifying the Cardan grille on transition statistics, and a content ablation test proving that real botanical content passed through the generation system.

We disclose that AI tools (Claude, Anthropic) were used for coding assistance, botanical annotation of manuscript pages, and systematic code review, as detailed in Section~2.7. All code, data, and annotations are publicly available at \texttt{github.com/leCheeseRoyale/voynich-lossy-generator}.

The corresponding author publishes under a pseudonym. Legal identity is available to the editor upon request.

\closing{Respectfully,}
\end{letter}
\end{document}
